\documentclass[reqno, answers]{exam}

\usepackage{amssymb, amsmath, amsthm, stmaryrd}
\usepackage{fullpage, parskip}
\usepackage{enumitem}
\usepackage{tikz}
\usepackage{ulem}
\usepackage{hyperref}

\title{Introduction to Computational Math Spring 2026 \\ Homework 04}
\date{Due: Friday, Feb 20, 2026, 11:59 PM}
\author{}

\begin{document}
\maketitle
\thispagestyle{empty}

Textbook = Driscoll and Braun (\url{https://fncbook.com/})

\begin{questions}

\question Textbook Exercise 2.7.3

Prove that for all vectors $\mathbf{x} \in \mathbb{R}^n$:
\begin{parts}
\part $\|\mathbf{x}\|_\infty \leq \|\mathbf{x}\|_2$
\part $\|\mathbf{x}\|_2 \leq \|\mathbf{x}\|_1$
\end{parts}

\begin{solution}
\textbf{(a) Prove $\|\mathbf{x}\|_\infty \leq \|\mathbf{x}\|_2$:}

Let $\mathbf{x} = (x_1, x_2, \ldots, x_n)^T \in \mathbb{R}^n$.

By definition:
\begin{itemize}
    \item $\|\mathbf{x}\|_\infty = \max_{1 \leq i \leq n} |x_i|$
    \item $\|\mathbf{x}\|_2 = \sqrt{\sum_{i=1}^n x_i^2}$
\end{itemize}

Let $k$ be the index where the maximum is attained, so $\|\mathbf{x}\|_\infty = |x_k|$.

Then:
\[
\|\mathbf{x}\|_2^2 = \sum_{i=1}^n x_i^2 \geq x_k^2 = |x_k|^2 = \|\mathbf{x}\|_\infty^2
\]

Taking square roots (both sides are non-negative):
\[
\|\mathbf{x}\|_2 \geq \|\mathbf{x}\|_\infty
\]

\textbf{(b) Prove $\|\mathbf{x}\|_2 \leq \|\mathbf{x}\|_1$:}

By definition:
\begin{itemize}
    \item $\|\mathbf{x}\|_1 = \sum_{i=1}^n |x_i|$
    \item $\|\mathbf{x}\|_2 = \sqrt{\sum_{i=1}^n x_i^2}$
\end{itemize}

We need to show $\|\mathbf{x}\|_2^2 \leq \|\mathbf{x}\|_1^2$, i.e.:
\[
\sum_{i=1}^n x_i^2 \leq \left(\sum_{i=1}^n |x_i|\right)^2
\]

Expanding the right side:
\[
\left(\sum_{i=1}^n |x_i|\right)^2 = \sum_{i=1}^n x_i^2 + 2\sum_{i < j} |x_i||x_j|
\]

Since $\sum_{i < j} |x_i||x_j| \geq 0$, we have:
\[
\sum_{i=1}^n x_i^2 \leq \sum_{i=1}^n x_i^2 + 2\sum_{i < j} |x_i||x_j| = \|\mathbf{x}\|_1^2
\]

Taking square roots: $\|\mathbf{x}\|_2 \leq \|\mathbf{x}\|_1$.
\end{solution}

\question Textbook Exercise 2.7.6

Let $A = \begin{bmatrix} -1 & 1 \\ 2 & 2 \end{bmatrix}$.
\begin{parts}
\part Find all vectors satisfying $\|\mathbf{x}\|_\infty = 1$ and $\|A\mathbf{x}\|_\infty = \|A\|_\infty$.
\part Find a vector satisfying $\|\mathbf{x}\|_1 = 1$ and $\|A\mathbf{x}\|_1 = \|A\|_1$.
\part Find a vector satisfying $\|\mathbf{x}\|_2 = 1$ such that $\|A\mathbf{x}\|_2 = \|A\|_2$. (Hint: Use the definition of $\|A\|_2$ as the maximum of $\|A\mathbf{x}\|_2$ over unit vectors, and parameterize the unit vectors as $(\cos\theta, \sin\theta)$. Then $\|A\mathbf{x}\|_2^2$ is a function of $\theta$ that is easy to maximize.)
\end{parts}

\begin{solution}
\textbf{(a) Finding vectors with $\|\mathbf{x}\|_\infty = 1$ and $\|A\mathbf{x}\|_\infty = \|A\|_\infty$:}

First, compute $\|A\|_\infty$ (maximum absolute row sum):
\begin{itemize}
    \item Row 1: $|-1| + |1| = 2$
    \item Row 2: $|2| + |2| = 4$
\end{itemize}

So $\|A\|_\infty = 4$.

For $\mathbf{x} = \begin{bmatrix} x_1 \\ x_2 \end{bmatrix}$ with $\|\mathbf{x}\|_\infty = 1$:
\[
A\mathbf{x} = \begin{bmatrix} -x_1 + x_2 \\ 2x_1 + 2x_2 \end{bmatrix}
\]

We need $\|A\mathbf{x}\|_\infty = 4$. This requires $|2x_1 + 2x_2| = 4$, i.e., $|x_1 + x_2| = 2$.

With $\|\mathbf{x}\|_\infty = 1$ (meaning $|x_1| \leq 1$ and $|x_2| \leq 1$, with at least one equality), the maximum of $|x_1 + x_2|$ is 2, achieved when $x_1 = x_2 = 1$ or $x_1 = x_2 = -1$.

\textbf{Answer:} $\mathbf{x} = \begin{bmatrix} 1 \\ 1 \end{bmatrix}$ and $\mathbf{x} = \begin{bmatrix} -1 \\ -1 \end{bmatrix}$.

\textbf{(b) Finding a vector with $\|\mathbf{x}\|_1 = 1$ and $\|A\mathbf{x}\|_1 = \|A\|_1$:}

First, compute $\|A\|_1$ (maximum absolute column sum):
\begin{itemize}
    \item Column 1: $|-1| + |2| = 3$
    \item Column 2: $|1| + |2| = 3$
\end{itemize}

So $\|A\|_1 = 3$.

For unit vectors in the 1-norm, choose $\mathbf{x} = \mathbf{e}_1 = \begin{bmatrix} 1 \\ 0 \end{bmatrix}$:
\[
A\mathbf{e}_1 = \begin{bmatrix} -1 \\ 2 \end{bmatrix}, \quad \|A\mathbf{e}_1\|_1 = |-1| + |2| = 3 = \|A\|_1
\]

\textbf{Answer:} $\mathbf{x} = \begin{bmatrix} 1 \\ 0 \end{bmatrix}$ (or $\begin{bmatrix} 0 \\ 1 \end{bmatrix}$).

\textbf{(c) Finding a vector with $\|\mathbf{x}\|_2 = 1$ and $\|A\mathbf{x}\|_2 = \|A\|_2$:}

Parameterize unit vectors as $\mathbf{x} = \begin{bmatrix} \cos\theta \\ \sin\theta \end{bmatrix}$.

\[
A\mathbf{x} = \begin{bmatrix} -\cos\theta + \sin\theta \\ 2\cos\theta + 2\sin\theta \end{bmatrix}
\]

\begin{align*}
\|A\mathbf{x}\|_2^2 &= (-\cos\theta + \sin\theta)^2 + (2\cos\theta + 2\sin\theta)^2 \\
&= \cos^2\theta - 2\sin\theta\cos\theta + \sin^2\theta + 4(\cos\theta + \sin\theta)^2 \\
&= 1 - 2\sin\theta\cos\theta + 4(\cos^2\theta + 2\sin\theta\cos\theta + \sin^2\theta) \\
&= 1 - \sin(2\theta) + 4(1 + \sin(2\theta)) \\
&= 5 + 3\sin(2\theta)
\end{align*}

This is maximized when $\sin(2\theta) = 1$, i.e., $2\theta = \frac{\pi}{2}$, so $\theta = \frac{\pi}{4}$ or $\theta = \frac{5\pi}{4}$.

Thus $\|A\|_2 = \sqrt{5 + 3} = \sqrt{8} = 2\sqrt{2}$.

\textbf{Answer:} $\mathbf{x} = \begin{bmatrix} \cos(\pi/4) \\ \sin(\pi/4) \end{bmatrix} = \begin{bmatrix} \frac{\sqrt{2}}{2} \\ \frac{\sqrt{2}}{2} \end{bmatrix} = \frac{\pm 1}{\sqrt{2}}\begin{bmatrix} 1 \\ 1 \end{bmatrix}$.
\end{solution}

\question Textbook Exercise 2.8.6

Define $A_n$ as the $n \times n$ matrix
\[
A_n = \begin{bmatrix} 1 & -2 & & \\ & 1 & -2 & \\ & & \ddots & \ddots \\ & & & 1 \end{bmatrix}.
\]
\begin{parts}
\part Write out $A_2^{-1}$ and $A_3^{-1}$.
\part Write out $A_n^{-1}$ in the general case $n > 1$. Make a clear argument why it is correct.
\part Using the $\infty$-norm, find $\kappa(A_n)$.
\end{parts}

\begin{solution}
\textbf{(a) Computing $A_2^{-1}$ and $A_3^{-1}$:}

For $n = 2$: $A_2 = \begin{bmatrix} 1 & -2 \\ 0 & 1 \end{bmatrix}$

Using the formula for $2 \times 2$ inverse:
\[
A_2^{-1} = \begin{bmatrix} 1 & 2 \\ 0 & 1 \end{bmatrix}
\]

For $n = 3$: $A_3 = \begin{bmatrix} 1 & -2 & 0 \\ 0 & 1 & -2 \\ 0 & 0 & 1 \end{bmatrix}$

We solve $A_3 \mathbf{x} = \mathbf{e}_i$ for each standard basis vector:
\[
A_3^{-1} = \begin{bmatrix} 1 & 2 & 4 \\ 0 & 1 & 2 \\ 0 & 0 & 1 \end{bmatrix}
\]

\textbf{(b) General formula for $A_n^{-1}$:}

The pattern shows that $A_n^{-1}$ is an upper triangular matrix with:
\[
(A_n^{-1})_{ij} = \begin{cases} 2^{j-i} & \text{if } i \leq j \\ 0 & \text{if } i > j \end{cases}
\]

\textbf{Verification:} Let $B = A_n^{-1}$ with $B_{ij} = 2^{j-i}$ for $i \leq j$. We verify $A_n B = I$:

For $(A_n B)_{ij}$:
\begin{itemize}
    \item If $i = 1$: $(A_n B)_{1j} = 1 \cdot B_{1j} = \delta_{1j}$ \checkmark
    \item If $i > 1$: $(A_n B)_{ij} = -2 \cdot B_{i-1,j} + 1 \cdot B_{ij}$
    \begin{itemize}
        \item If $j < i$: Both terms are 0, so $(A_n B)_{ij} = 0$ \checkmark
        \item If $j > i$: $2^{j-i} - 2 \cdot 2^{j-i-1} = 2^{j-i} - 2^{j-i} = 0$ \checkmark
        \item If $j = i$: $1 - 0 = 1$ \checkmark
    \end{itemize}
\end{itemize}

\textbf{(c) Finding $\kappa(A_n)$ in the $\infty$-norm:}

$\|A_n\|_\infty$ (max absolute row sum):
\begin{itemize}
    \item Row $1$ to $n-1$: $|1| + |-2| = 3$
    \item Row $n$: $|1| = 1$
\end{itemize}

So $\|A_n\|_\infty = 3$ for $n \geq 2$.

$\|A_n^{-1}\|_\infty$ (max absolute row sum of $A_n^{-1}$):

Row $i$ sum: $\sum_{j=i}^{n} 2^{j-i} = 2^0 + 2^1 + \cdots + 2^{n-i} = 2^{n-i+1} - 1$

Maximum at row $1$: $\|A_n^{-1}\|_\infty = 2^n - 1$

\textbf{Answer:} 
\[
\boxed{\kappa(A_n) = \|A_n\|_\infty \cdot \|A_n^{-1}\|_\infty = 3(2^n - 1)}
\]
\end{solution}

\question Textbook Exercise 2.8.8

Let $D$ be a diagonal $n \times n$ matrix, not necessarily invertible. Prove that in the 1-norm,
\[
\kappa(D) = \frac{\max_i |D_{ii}|}{\min_i |D_{ii}|}
\]
(Hint: See Exercise 2.7.10.)

\begin{solution}
Let $D = \text{diag}(d_1, d_2, \ldots, d_n)$ where $d_i = D_{ii}$.

\textbf{Case 1: $D$ is invertible (all $d_i \neq 0$)}

For a diagonal matrix, $D^{-1} = \text{diag}(1/d_1, 1/d_2, \ldots, 1/d_n)$.

For the 1-norm (maximum absolute column sum), since $D$ is diagonal:
\[
\|D\|_1 = \max_j \sum_i |D_{ij}| = \max_j |d_j| = \max_i |d_i|
\]

Similarly:
\[
\|D^{-1}\|_1 = \max_i |1/d_i| = \frac{1}{\min_i |d_i|}
\]

Therefore:
\[
\kappa(D) = \|D\|_1 \cdot \|D^{-1}\|_1 = \max_i |d_i| \cdot \frac{1}{\min_i |d_i|} = \frac{\max_i |D_{ii}|}{\min_i |D_{ii}|}
\]

\textbf{Case 2: $D$ is singular (some $d_i = 0$)}

If $\min_i |d_i| = 0$, then $D$ is singular and $\kappa(D) = \infty$ by convention.

The formula $\frac{\max_i |D_{ii}|}{\min_i |D_{ii}|}$ also gives $\infty$ when $\min_i |D_{ii}| = 0$ (assuming at least one $d_i \neq 0$).

Thus, the formula holds in both cases:
\[
\boxed{\kappa(D) = \frac{\max_i |D_{ii}|}{\min_i |D_{ii}|}}
\]
\end{solution}

\end{questions}

\end{document}